\section{評価の枠組み}
\label{sec_eval}
本研究では，手法やラベルの定義の改訂に合わせて評価用データを拡充してきた．
本章では，最終的に用いた3種類の評価用データと評価指標を定義し，それぞれを導入した理由を述べる．

\subsection{目視ラベルの付与}
評価に用いる正解のラベルは，判断者が画像を1枚ずつ見て付与した．
以下では，このラベルを目視ラベルと呼ぶ．
目視ラベルの付与は，画像を表示して内容を確認できる AI エージェントである Claude Code が，その時点の定義書に従って行った．
付与の際は，画像を20枚ずつ並べた一覧画像を作成し，各画像に番号とお題 ID を表示した．
評価する手法の予測は一覧画像に表示せず，予測に判断が引きずられないようにした．
目視ラベルは全て1名の判断者によるものであり，判断者間の一致度は測っていない．
この点は本研究の限界であり，\ref{sec_discussion}章で改めて述べる．

\subsection{通常の評価用データ}
通常の評価用データは，全てのお題から無作為に抽出した440枚と，少ない種類を補うためにメタデータで絞り込んだ200枚の計640枚からなる．
メタデータによる絞り込みでは，\texttt{type} 列が穴埋め，文字，イラストのお題と，カルタ，漫画，アニメのタグが付いたお題から抽出した．
最初の実験では440枚であり，2回目の実験で無作為抽出の200枚を加えて640枚とした．
目視ラベルは，ラベルの定義を改訂するたびに付け直した．

プロンプトの改訂を評価用データの誤りを見ながら行うと，評価用データでの精度は楽観的になる．
そこで，3軸のラベルを導入した段階から，640枚を形式で層化して320枚ずつに二分した．
以下では，プロンプトや構成の選択に用いる半分を調整用データ（dev），最終的な精度の報告にのみ用いる半分を最終評価用データ（test）と呼ぶ．
また，生徒のハイパーパラメータは，評価用データではなく，教師ラベルの一部を取り分けた検証用分割で選んだ．

\subsection{少ない区分の評価用データ}
通常の評価用データには，出現率の低いクラスがほとんど含まれない．
例えば，絵画は640枚中5枚，CG は4枚であり，1枚の正誤で F1 が大きく変わる．
そこで，ゲーム，CG，絵画，カルタの4区分を多く含む評価用データを別に作成した．
以下では，出現率が数\%以下のこれらのクラスを少ない区分と呼ぶ．

第1版は，生徒が各区分である確率の上位2,000枚から，区分ごとに50枚を無作為に抽出した200枚である．
第2版では，学習にも評価にも用いていない294,802枚の画像に生徒の確率を付け，区分 $ c $ ごとに確率が0.5以上の層と，0.05以上0.5未満の層から60枚ずつを抽出した．
さらに全体から無作為に120枚を加え，計600枚とした．
第1版と第2版を合わせた800枚を，抽出の層ごとに dev と test に二分した．
この評価用データは確率で絞り込んだ画像からなるため，区分の出現率は実際より高い．
したがって，このデータでの値は手法間の比較にのみ用いる．

\subsection{実データの目視ラベル}
通常の評価用データには，ボケての典型的な画像の一部が少ないことが，研究の後半で明らかになった．
そこで，学習にも評価用データにも用いていない実データから，2通りの方法で画像を抽出し，目視ラベルを付けた．
1つ目は無作為抽出の160枚であり，実際の分布での精度を直接表す．
2つ目は層別抽出の220枚である．
層別抽出では，生徒が予測したクラスのうち実データで最も出現率の低いものを各画像の層とし，11の層から20枚ずつを抽出した．
これにより，絵画，カルタ，ゲーム，お題文などの少ないクラスを各15枚から20枚程度確保した．
無作為抽出と層別抽出を合わせて，実データの目視ラベルと呼ぶ．

層別抽出の画像は抽出の確率が層によって異なるため，単純に集計すると少ないクラスの影響が過大になる．
そこで，各画像に抽出確率の逆数の重みを付けて，実際の分布での精度を推定した．
画像 $ i $ の重み $ w_i $ と，ある指標の推定値 $ \hat{\mu} $ を(\ref{eq_ipw})に示す．

\begin{equation}
    w_i = \frac{N_{h(i)}}{n_{h(i)}},
    \qquad
    \hat{\mu} = \frac{\sum_{i} w_i z_i}{\sum_{i} w_i}
    \label{eq_ipw}
\end{equation}

\noindent
ここで $ h(i) $ が画像 $ i $ の層，$ N_h $ が実データ294,802枚のうち層 $ h $ に属する画像の数，$ n_h $ が層 $ h $ から抽出した画像の数，$ z_i \in \{0, 1\} $ が画像 $ i $ で指標の条件が成り立つかを表す値である．
無作為抽出の160枚は，全体の画像数を抽出数で割った重みで同じ式に加えた．
以下では，この値を推定値と呼ぶ．

\subsection{評価指標}
形式のように1つのクラスを選ぶ軸は，正解率とマクロ F1 で評価した．
マクロ F1 を(\ref{eq_macro_f1})に示す．

\begin{equation}
    {\rm F1}_{\rm macro} = \frac{1}{K} \sum_{k=1}^{K} \frac{2 P_k R_k}{P_k + R_k}
    \label{eq_macro_f1}
\end{equation}

\noindent
ここで $ K $ がクラス数，$ P_k $ と $ R_k $ がクラス $ k $ の適合率と再現率である．
マクロ F1 は，件数の少ないクラスの性能を件数の多いクラスと等しく反映する．
属性と媒体のように複数の項目を選びうる軸は，項目ごとに2値分類の F1 を求め，評価用データに正例がある項目の F1 を平均して，その軸のマクロ F1 とした．
総合の指標として，形式，属性，媒体の3軸のマクロ F1 の平均を用い，以下では3軸の平均と呼ぶ．

件数の少ない項目は F1 の推定が不安定であるため，正例が10件以上ある項目のみで求めた平均も併用し，頑健な指標と呼ぶ．
少ない区分の評価用データでは，各区分の F1 の平均を用いた．
実データの目視ラベルでは，全ての項目が目視ラベルと一致した画像の割合を全項目一致と呼び，(\ref{eq_ipw})による推定値を主な指標とした．
2つの構成の全項目一致の差の信頼区間は，画像を単位とする対応のある自助法により，2,000回の再標本化で求めた．
